\chapter{Evaluation}
\label{ch:evaluation}

% Budget and sources: docs/thesis/outline.md. Every claim about behaviour points at a test, a figure or a number in data/.


Every number in this chapter comes from run \code{2026-10-19\_1200} in \code{docs/thesis/data}, produced by
\code{make thesis:eval} on a single laptop (Intel i9-14900HX, 32 cores, 30~GB; Go~1.27.1, PostgreSQL~18.6) with
the whole stack running locally. The run has two parts. The first executes the service's test suites, unit and
database-backed together, and records what they assert; the second drives the running service with a paced load
and reads its own metrics back for exactly that window. Section~\ref{sec:eval-decisions} reports a separate,
earlier measurement of the decision model.

\section{Correctness across regimes}
\label{sec:eval-correctness}

The run's suites pass in full: 189 tests across 19 packages, none failing. The ones that bear on the central claim
of chapter~\ref{ch:design}, that one state machine carries four regimes, are summarised in
figure~\ref{fig:eval-coverage}. For each regime the test explores every state reachable from \emph{initiated},
tries every event in each, and records which pairs are accepted and which refused. It then fails unless each
accepted transition is walked by at least one named scenario, a dispute taken from \emph{initiated} to
\emph{closed} along a path a reader can follow, such as ``Regulation~E: chargeback lost without evidence, credit
reversed''.

\begin{figure}[htbp]
\centering
\small
\begin{tabular}{lrrrrr}
\toprule
Regime & States & Accepted & Walked & Scenarios & Refused pairs \\
\midrule
SEPA direct debit & 6 & 10 & 10 & 6 & 62 \\
PSD2 card & 11 & 18 & 18 & 8 & 114 \\
Regulation~E & 13 & 20 & 20 & 8 & 136 \\
Regulation~Z & 10 & 16 & 16 & 7 & 104 \\
\midrule
Total & & 64 & 64 & 29 & 416 \\
\bottomrule
\end{tabular}
\caption{Transition coverage per regime in run \code{2026-10-19\_1200} (\code{correctness.json},
\code{TestTransitionCoverage}). States are those reachable from \emph{initiated}; accepted and refused count the
pairs of reachable state and event; walked counts the accepted transitions some named scenario takes.}
\label{fig:eval-coverage}
\end{figure}

Three other tests close the remaining gaps in the argument. The set of reachable states of each regime is pinned,
so that adding or removing a state is a deliberate change with its own diff rather than an accident. Every
reachable state must be able to reach \emph{closed}, so no regime can strand a dispute. And the events the
engine offers in each reachable state must be exactly the events it would accept, which is what lets the
workbench show only the actions that will succeed. The refusals that carry the regulation, such as a SEPA claim
that tries to file a chargeback or a Regulation~Z dispute that tries to issue a refund, are named tests of their
own. Before the coverage test existed, the named scenarios walked only 28 of the 64 accepted transitions; the
other 36 were covered only by the consistency check, which shows that the engine agrees with itself but not that
a path makes sense. The 23 scenarios added to close that gap are now part of the suite, and removing a scenario
that is the only one to walk some transition makes the build fail, naming that transition.

\section{The log, concurrency and retries}
\label{sec:eval-log}

Section~\ref{sec:design-state-machine} relies on the dispute's history being trustworthy, and
section~\ref{sec:design-core} relies on retries being safe. Four tests against a real PostgreSQL database
establish both. The first tries to rewrite history directly in the database, with an update of an event's actor,
a delete and a truncate of the event table, and each is refused by the database itself; the dispute still has
exactly the one event it started with. The second starts eight writers at the same instant, each applying the
same event to the same dispute: exactly one succeeds, the other seven receive a conflict, and the dispute ends at
version two with two events, so the race neither lost a write nor applied one twice. The third repeats a request
with the same idempotency key on a different connection and receives the stored response rather than a second
dispute, while the same key with a different body is refused. The fourth applies a transition the regime forbids
and checks that nothing at all was written, neither an event nor a version.

The load run exercises the same mechanisms under traffic. In two minutes at about twenty requests a second it
opened 479 disputes and applied 958 transitions, replayed one creation in ten with its original key, and
deliberately attempted one forbidden event per dispute. The response counts account for every request: 526
creations answered \code{201} (479 new and 47 replays), 1,437 reads and transitions answered \code{200}, and
the 479 forbidden events each answered \code{409}, with no server errors. The load generator checks only that a
replay is answered like the original; that a replay never creates a second dispute is the database test above.

\section{Isolation between tenants and teams}
\label{sec:eval-isolation}

The isolation claim of section~\ref{sec:design-tenancy} is that a tenant's data is invisible to every other
tenant even if the application code is wrong. It is therefore tested below the application, against the
database role the service actually runs as. Under that role, a second tenant cannot read the first tenant's
dispute, cannot apply an event to it, and cannot open a dispute on its transaction; in all three cases the answer
is ``not found'', so the attempt does not even learn that the dispute exists. The same idempotency key used by two
tenants is two separate requests, not a replay of the other's response. A query on a connection that has not
been bound to a tenant sees no rows at all, so a forgotten binding fails closed. And an attempt to insert a row
carrying another tenant's identifier is rejected by the policy with PostgreSQL's insufficient-privilege error,
so a row cannot be written into another tenant either.

Team scope, the second line of section~\ref{sec:design-teams}, is tested the same way. An analyst who belongs
only to another team cannot read, list or advance a dispute, and a raw query in that analyst's context sees none
of its events. The test also checks the property that made the policies restrictive rather than permissive: a team
lead of another tenant still sees nothing, so team scope can only narrow what tenant isolation allows. The authorization
policies on top are checked separately, over generated sets of grants, for the three properties that matter:
no set lets an investigator pay out their own dispute, reach a dispute outside their teams, or administer the
tenant.

\section{Latency under load}
\label{sec:eval-latency}

The load run measures what a single instance costs per request, not what a deployment could carry, since the
whole stack shared one machine. It paced four lifecycle iterations a second, about twenty requests, for two
minutes against one tenant, and read the server's own latency histograms back from Prometheus for that window
(figure~\ref{fig:eval-latency}).

\begin{figure}[htbp]
\centering
\small
\begin{tabular}{lrrr}
\toprule
Route & p50 & p95 & p99 \\
\midrule
Read a dispute (\code{GET /disputes/\{id\}}) & 3.3 & 7.2 & 9.5 \\
Apply an event (\code{POST /disputes/\{id\}/events}) & 5.1 & 9.6 & 16.8 \\
Open a dispute (\code{POST /disputes}) & 13.7 & 24.5 & 42.7 \\
\midrule
All requests, as the client saw them & 5.1 & 15.6 & 20.9 \\
\bottomrule
\end{tabular}
\caption{Latency in milliseconds per route, server side, in run \code{2026-10-19\_1200} (\code{load.json}),
at about 20 requests a second for 120 seconds.}
\label{fig:eval-latency}
\end{figure}

The numbers have the shape the design predicts. A read is a handful of queries under row-level security and stays
under ten milliseconds at the 99th percentile. An event costs more because it does more inside one transaction:
the authorization decision, the state machine, the clocks, the ledger postings, the notices it owes and, where
money moves, the call to the simulated core. Opening a dispute costs most, because it also starts the clocks, scores the dispute and
writes the first notice. PostgreSQL itself was never the bottleneck: its 95th percentile per query stayed under
one millisecond (0.97~ms), and the service never ran more than 61 goroutines or held more than 25~MB of heap.

Two comparisons put these numbers in context. Measured the same way on the same machine before authorization
existed, the median read took 2.1 milliseconds and the median event 3.0; adding a policy decision per operation
and team scope in the database raised them to 4.2 and 5.0 milliseconds (runs \code{2026-09-30\_1050} and
\code{2026-09-30\_1057}). The difference is round trips rather than policy evaluation, which costs about 25
microseconds a decision. Against that later run, this one's reads and events are no slower, while
the 99th percentile of opening a dispute rose from 25 to 43 milliseconds although the service's code is the same.
This run is the first whose database-backed suites ran against the same database immediately before the load,
which is the likely cause; it has not been isolated, and the figure is reported as measured.

\section{Are the model's proposals good enough to show?}
\label{sec:eval-decisions}

The feature of section~\ref{sec:impl-decisions} is only defensible if the proposals are usually
right, so it was measured before it was built. The labelled set is drawn from this system's own
vocabulary: the four reasons, the yes and no questions of the unauthorised questionnaire, and the
fifteen lifecycle states. It holds 72 complaints, 113 questionnaire answers across 40 replies, and
69 filter fields across 50 queries. Every part carries a deliberately hard band, and roughly a
third is Hungarian, since both seeded tenants are Hungarian banks and a customer writes in the
language they complain in. Two models were run against it: an open one on a consumer GPU, and a
hosted one (table~\ref{fig:decisions}).

\begin{figure}[htbp]
\centering
\small
\begin{tabular}{lrrrrr}
\toprule
Model & Reason & Questionnaire & Search & ECE & One question \\
\midrule
Open, English checkpoint & 0.597 & 0.469 & 0.812 & 0.401 & 17 ms \\
Open, multilingual & 0.639 & 0.522 & 0.623 & 0.213 & 8 ms \\
Open, typed decisions & 0.639 & 0.611 & 0.696 & 0.595 & 16 ms \\
Hosted & \textbf{0.986} & \textbf{0.973} & \textbf{0.928} & \textbf{0.011} & 574 ms \\
\midrule
Chance & 0.25 & 0.50 & & & \\
\bottomrule
\end{tabular}
\caption{Accuracy on the labelled set (\code{eval/decisions}), expected calibration error on the
reason task, and the latency of one question. The open model runs locally on an RTX~4070; the
hosted one is a network call costing \$0.0000216. Chance is one in four for the reason task and
one in two for the questionnaire.}
\label{fig:decisions}
\end{figure}

The open model is not accurate enough for any of the three uses, and the reason it fails is not
only accuracy. Its expected calibration error ranges from 0.21 to 0.60, and it warns on loading
that its temperatures fall outside the calibrated range, so its confidence cannot be used to decide
when to stay quiet. The hosted model's error of 0.011 is what makes the confidence gate meaningful:
the single reason it got wrong, a rental company adding charges the customer never agreed to, came
back at 0.50 with the probability split 0.63 to 0.37 between two readings a person would also
argue about. Its three questionnaire mistakes were all the same question, which asked whether the
card was in the customer's possession and so conflated owning it with carrying it; the question was
reworded rather than the model replaced.

The honest limit of this measurement is that the set is synthetic. No corpus of real customer
language exists for this system, so the prose was written rather than collected, and it was
labelled by its author. It therefore measures how well a model recovers an intent that was
deliberately encoded, which is an easier task than reading what a real customer writes. The number
the running system reports, the share of proposals an analyst keeps, is the one that will answer
the question properly, and it is recorded from the first dispute opened this way.

\section{Limitations}
\label{sec:limitations}

The evaluation shows the shape of the system's behaviour, not its capacity. The load ran on one machine, with the
service, the database, the identity provider and the telemetry stack sharing it and talking over loopback, from
one tenant, with synthetic disputes whose mix was chosen rather than observed. A deployment would add network
latency between the parts and contention between tenants that this run cannot show, and the numbers say nothing
about how many disputes a bank of a given size would produce.

The correctness results are tests, not proofs. Coverage is measured over transitions, not over paths, so a fault
that appears only after a particular sequence of steps would not be caught merely because each step is walked
somewhere. The regimes themselves are modelled as chapter~\ref{ch:requirements} describes and no further: the
statutory extensions of the deadlines are not modelled, and every dispute runs the standard clocks; the
contested rounds of a chargeback after the merchant's first answer are folded into one evidence step; and the
banking core and the card network are simulated, so the behaviour of a real core under failure is assumed from
its interface rather than observed. The analysts' identity is delegated to the identity provider, whose own
protections, such as multi-factor authentication, are configured rather than evaluated here. Finally, the decision
model was measured on a synthetic, self-labelled set (section~\ref{sec:eval-decisions}); the share of its proposals
analysts actually keep is recorded by the running system and is the measurement that will settle its value.
